\documentclass{article}

\usepackage{project_440_550}
% Please submit it as is here, with line numbers.
% If you'd like a "less draft"-looking version for your website or something after:
%     \usepackage[final]{project_440_550}   % keeps the footer but kills line numbers,
%     \usepackage[preprint]{project_440_550}   % removes both


\usepackage[utf8]{inputenc} % allow utf-8 input
\usepackage[T1]{fontenc}    % use 8-bit T1 fonts

\usepackage[USenglish]{babel}  % there's a "canadian" option, but it's an alias for USenglish,
                               % and for some reason it makes csqoutes behave differently...
\usepackage{csquotes}       % smarter handling of quotes (used by biblatex)

\usepackage{booktabs}       % professional-quality tables
\usepackage{amsfonts}       % blackboard math symbols
\usepackage{nicefrac}       % compact symbols for 1/2, etc.
\usepackage{microtype}      % microtypography
\usepackage{xcolor}         % colors
\usepackage{hyperref}       % hyperlinks
\usepackage{url}            % simple URL typesetting

\usepackage{multirow}

% I recommend the biblatex package.
% If you hate it for some reason, though, you can use natbib instead:
% comment out this block and uncomment the next one.

\usepackage[style=authoryear,maxbibnames=30,natbib]{biblatex}
\addbibresource{refs.bib}
\renewbibmacro{in:}{}  % drops a silly "In:" from biblatex format
\DeclareDelimFormat{nameyeardelim}{\addspace} % remove a comma i dislike, doesn't matter
\DeclareNameAlias{sortname}{given-family}  % avoid kinda-weird bib printing order

%\usepackage[round]{natbib}
%\newcommand{\printbibliography}{\bibliographystyle{plainnat}\bibliography{refs}}

\newcommand{\rc}[2]{%
  \textcolor{red}{\fbox{\textbf{#1: #2}}}%
}

\newcommand{\hs}[1]{\textcolor{blue}{\textbf{[HS: #1]}}}
\newcommand{\ek}[1]{\textcolor{green}{\textbf{[EK: #1]}}}
\newcommand{\ag}[1]{\textcolor{orange}{\textbf{[AG: #1]}}}


\usepackage[capitalize,noabbrev]{cleveref}


\title{KAN it learn? Interpretable N-Body Physics}


% The \author macro works with any number of authors. There are two commands
% used to separate the names and addresses of multiple authors: \And and \AND.
%
% Using \And between authors leaves it to LaTeX to determine where to break the
% lines. Using \AND forces a line break at that point. So, if LaTeX puts 3 of 4
% authors names on the first line, and the last on the second line, try using
% \AND instead of \And before the third author name.

\author{%
  Hanson Sun\\
  \texttt{jsun33@students.cs.ubc.ca}
  \And
  Eitan Klinger\\
  \texttt{ekling01@students.cs.ubc.ca}
  \And
  Arthur Guo\\
  \texttt{aguo07@students.cs.ubc.ca}
}

\begin{document}
\maketitle


\begin{abstract}
    Our project literally sucks and doesnt work
\end{abstract}

\section{Introduction}

Many scientific systems can be modeled as interacting objects governed by underlying physical laws. A key goal in machine learning for science is not only to predict these systems accurately, but also to recover interpretable structure from the learned model.

We study the $n$-body problem as a testbed for this goal. Given the initial states of several bodies, the task is to predict their future trajectories under pairwise forces. Prior work has shown that graph neural networks (GNNs) are a strong fit for this setting because they naturally represent objects as nodes and interactions as edges. In particular, \cite{cranmer2020discoveringsymbolicmodelsdeep} showed that a GNN trained on simulated $n$-body dynamics can be combined with symbolic regression to recover analytical force laws.

A limitation of this approach is that the GNN itself is not directly interpretable. The learned representation is still a black box, and symbolic regression is a separate post-hoc step. Recent work on Kolmogorov-Arnold Networks (KANs) suggests a possible alternative. KANs replace standard MLP-style node activations with learnable univariate functions on edges, which can make the learned model easier to inspect. \cite{liu2024kankolmogorovarnoldnetworks} showed that KANs can be effective while remaining more transparent than standard MLPs.

In this project, we ask whether this idea can be extended to relational physical systems. We introduce a hybrid Graph KAN (GKAN) architecture, where KAN layers replace MLP components inside a message-passing GNN. We evaluate this model on simulated $n$-body dynamics and compare it against standard GNN baselines, with a focus on both prediction quality and interpretability.

\section{Related Work}

This project builds on three main lines of work: graph-based physics learning, symbolic discovery, and interpretable neural architectures.

\cite{cranmer2020discoveringsymbolicmodelsdeep} is the closest prior work. They trained GNNs on simulated $n$-body systems and used symbolic regression to recover physical laws from learned interaction functions. Their results show that graph models can discover useful structure, but interpretability comes only after a separate extraction step.

Other work has strengthened graph models for physical simulation. \cite{thangamuthu2023unravellingperformancephysicsinformedgraph} compared several physics-informed GNNs across multiple dynamical systems, while Dynami-CAL GraphNet \parencite{sharma2025dynamicalgraphnetphysicsinformedgraph} adds conservation and symmetry constraints directly into message passing. These methods improve performance, but do not directly address whether the learned interactions are easy to interpret.

KANs provide a different route to interpretability. \cite{liu2024kankolmogorovarnoldnetworks} showed that KANs can represent scientific relationships through learnable univariate functions, and KAN 2.0 adds multiplicative structure that is especially relevant for physical equations. Related continuous-time work such as KAN-ODEs \parencite{Koenig_2024} and Graph KAN-ODEs \parencite{cappi2025unveilingactualperformanceneuralbased} further motivate combining KANs with dynamical systems and graph structure.

\section{Methodology}

We evaluate GKANs on synthetic $n$-body simulation data generated with a leapfrog integrator. The main dataset uses Newtonian gravity, and we also generate additional force laws for generalization tests. Each example consists of initial positions, velocities, and masses, with targets given by the future trajectories.

\hs{we need to provide the formulas of the force laws we used somewhere, maybe in appendix}
\hs{we also should provide baseline specification somewhere, ie how many layers, width, etc}

We compare two models: a graph convolutional network baseline following Cranmer et al.\cite{cranmer2020discoveringsymbolicmodelsdeep}, and a GKAN that replaces selected MLP components with KAN layers. Both models are trained on the same data, with the same loss and optimization setup, to keep the comparison fair.

We evaluate models using trajectory prediction error on held-out simulations. We also examine interpretability through sparsity, visual structure of learned edge functions, and how easily the model can reflect physically meaningful relationships.

We organize our experiments around a few questions: whether GKANs can match GNN performance at similar parameter counts, whether sparse regularization makes learned functions easier to interpret, whether smaller KAN-style architectures are sufficient, whether physically motivated feature engineering helps, and whether the same trends hold under different force laws such as Hooke's law.

\hs{will probably skip multkans, dont think i got anything useful there. can include if theres space}

\subsection{Regularization}
Sparsification can be achieved in KAN layers through L1 regularization of weights, as in traditional MLP layers, though an additional entropy regularizer is often needed to make sufficiently sparse KANs \cite{liu2024kankolmogorovarnoldnetworks}. This suite of experiments serves to both validate regularization patterns found in traditional KANs, as well as find ideal regularization hyperparameters to encourage sparse and interpretable networks. To do so, we run a sweep through different magnitudes of L1 regularization, entropy regularization, and both combined, comparing the models to each other, as well as a control with no regularization. Doing so helps identify the efficacy of different regularizations in creating sparser models. 

Each GKAN network uses a node network of dimensions [7, 8, 2], a message network of dimensions [10, 12, 12, 2], and is trained over 50 epochs. Once all models are trained, they are compared in both training loss, validation loss against a test set, as well as visually compared for how sparse the networks are. The ideal regularization approach is determined by which hyperparameters provide a relatively sparse and interpretable network without sacrificing performance on the validation set.

Idea: try to test out and find the best regularization hyperparameters to encourage sparse networks. all experiments done using gravity, n\_train = 500, n\_val = 50, msg\_width =  msg\_width: [10, 12, 12, 2], node\_width: [7, 8, 2], over 50 epochs using adam
\begin{itemize}
    \item control with no regularization
    \item sweep of l1 regularization: 1e-3, 1e-4, 1e-5 \ek{maybe drop the 1e-5, since it's the only case using that number and isn't important}
    \item sweep of entropy regularization: 1e-3, 1e-4
    \item sparse initialization 
    \item sweep of both l1 and entropy regularization: 1e-4, 1e-3, 1e-1, 1e1, 1e2
\end{itemize}
\hs{justify choice of regularization (L1 and entropy)}

\subsection{How small KAN we get?}

The Kolmogorov-Arnold (KA) representation theorem states that any multivariate function can be decomposed into a finite number of univariate functions, combined via addition and function composition. \ag{Who do I cite here? K and A from the 50s? Also, re-stating this stuff isn't actually that useful...}

\[
f(x_{1 \ldots n}) = \sum_{i = 0}^{2n} \Phi_{i} \left( \sum_{j = 1}^{n} \phi_{i, j}(x_{j}) \right)
\]

This suggests the existence of models with significantly fewer parameters that can still represent the function we want to learn. In this section, we design and evaluate several smaller models to determine if the GKAN can perform well without extreme overparameterization.

To design a smaller model, we work backwards from a proposed ``feasible route'' of function compositions and additions that yields the desired force law. In other words, we set layer widths and depths based on the number of arithmetic steps needed to decompose the desired force law into univariate functions by hand. We then observe whether the model is capable of recovering the desired force law by ``undoing'' our decomposition.

To illustrate this process, Figure \ref{fig:feasible-route-small} is one example of a feasible route. In this example, the message network consists of 4 layers ($10 \to 6 \to 3 \to 4 \to 2$) because the acceleration vector can be computed in 4 arithmetic steps; the node update network consists of just 1 layer ($7 \to 2$) because it only needs to aggregate acceleration vectors from the messages. \ag{I suppose we should explain the other 3 experiments in an appendix.}

\begin{figure}[h]

\centering
\begin{tabular}{c|l}
input & $m_{i}$, $s_{x,i}$, $s_{y,i}$, $v_{x,i}$, $v_{y,i}$, $m_{j}$, $s_{x,j}$, $s_{y,j}$, $v_{x,j}$, $v_{y,j}$ \\
\hline
\multirow{6}{*}{1} & $h_{1,1} = s_{x,j} - s_{x,i} \propto \Delta x$ \\
& $h_{1,2} = s_{y,j} - s_{y,i} \propto \Delta y$ \\
& $h_{1,3} = m_{j} + s_{x,j} - s_{x,i} \propto m_{j} + \Delta x$ \\
& $h_{1,4} = m_{j} - s_{x,j} + s_{x,i} \propto m_{j} - \Delta x$ \\
& $h_{1,5} = m_{j} + s_{y,j} - s_{y,i} \propto m_{j} + \Delta y$ \\
& $h_{1,6} = m_{j} - s_{y,j} + s_{y,i} \propto m_{j} - \Delta y$ \\
\hline
\multirow{3}{*}{2} & $h_{2, 1} = h_{1, 1}^{2} + h_{1, 2}^{2} \propto \|r\|^{2}$ \\
& $h_{2, 2} = h_{1,3}^{2} - h_{1,4}^{2} \propto m_{j} \Delta x$ \\
& $h_{2, 3} = h_{1,5}^{2} - h_{1,6}^{2} \propto m_{j} \Delta y$ \\
\hline
\multirow{4}{*}{3} & $h_{3,1} = h_{2,1}^{-1.5} + h_{2,2} \propto \|r\|^{-3} + m_{j} \Delta x$ \\
& $h_{3,2} = h_{2,1}^{-1.5} - h_{2,2} \propto \|r\|^{-3} - m_{j} \Delta x$ \\
& $h_{3,3} = h_{2,1}^{-1.5} + h_{2,3} \propto \|r\|^{-3} + m_{j} \Delta y$ \\
& $h_{3,4} = h_{2,1}^{-1.5} - h_{2,3} \propto \|r\|^{-3} - m_{j} \Delta y$ \\
\hline
\multirow{2}{*}{4} & $h_{4,1} = h_{3,1}^{2} - h_{3,2}^{2} \propto \|r\|^{-3} m_{j} \Delta x$ \\
& $h_{4,2} = h_{3,3}^{2} - h_{3,4}^{2} \propto \|r\|^{-3} m_{j} \Delta y$ \\
\hline
output & $a_{x,i} = h_{4,1}$, $a_{y,i} = h_{4,2}$ \\
\end{tabular}
\caption{A feasible route. \ag{Need a better way to present this.}}
\label{fig:feasible-route-small}

\end{figure}

In this section, we also experiment with feature engineering by augmenting the dataset with computed features. Specifically, we are interested in computed features that allow us to ``skip steps'' during the decomposition process. For example, in the decomposition described by Figure \ref{fig:feasible-route-small}, it is necessary to compute the intermediate values $\Delta x$ and $\Delta y$ in the first layer. By providing these as inputs, we can eliminate 1 layer.

In total, we designed and evaluated 4 smaller models with several combinations of computed features as listed in Figure \ref{fig:experiments-small}. For comparison, the baseline GNN model based on \cite{cranmer2020discoveringsymbolicmodelsdeep} has 300 hidden MLP layers and 427\,002 parameters.

\begin{figure}[h]

\centering
\begin{tabular}{l|l|l|l}
experiment name & message network & parameters & computed features \\
\hline
\texttt{small} & 4 layers ($10 \to 6 \to 3 \to 4 \to 2$) & 2220 & \\
\texttt{small\_rel\_pos} & 3 layers ($12 \to 5 \to 4 \to 2$) & 2020 & $\Delta x$, $\Delta y$ \\
\texttt{small\_rel\_dsq} & 3 layers ($13 \to 5 \to 4 \to 2$) & 2112 & $\Delta x$, $\Delta y$, $\|r\|^{+2}$\\
\texttt{small\_rel\_idc} & 3 layers ($13 \to 4 \to 4 \to 2$) & 1824 & $\Delta x$, $\Delta y, \|r\|^{-3}$ \\
\end{tabular}
\caption{The smaller models.}
\label{fig:experiments-small}

\end{figure}

The node update network always consists of 1 layer ($7 \to 2$). The models are all trained for 50 epochs on the same dataset, which consists of 200 training trajectories and 20 validation trajectories. For exact details on the data generation process and training regime, as well as an outline of the feasible route for each model, refer to the appendices. \ag{Use a ref here.}

\section{Results}
\hs{maybe have one big table summarizing losses for different configs (include maximal model somehow?)}

\subsection{Regularization}
\ek{insert table with training loss, val loss at 50 epochs? TODO: load respective checkpoints to get actual numbers, and plot the losses on a graph. make a script so we can do this for other models}


I think most of this actually goes in discussion?
\begin{itemize}
    \item all ts sucks ass!!!!
    \item regularization did not encourage sparsity
    \item just made model learn worse as regularization increased (regularization penalty overwhelming loss function?)
    \item like no discernible diff :(
    \item sparse init just sucked \ek{barely worth mentioning}
    \item regularization combined 1e1, 1e2 literally did not learn like the training error stayed flat. same for sparse init
\end{itemize}
TODO: generate loss plots overlaid, regenerate visualizations for experiments that are missing them

\subsection{How small KAN we get?}

Recall that we designed and evaluated 4 smaller models. The training loss and validation loss are summarized in Figure \ref{fig:results-small}.

\begin{figure}[h]

\centering
\begin{tabular}{l|c|c}
experiment name & training loss & validation loss \\
\hline
\texttt{baseline} & $1.592 \times 10^{-2}$ & $1.281 \times 10^{-2}$ \\
\texttt{small} & $1.130 \times 10^{-1}$ & $1.185 \times 10^{-1}$ \\
\texttt{small\_rel\_pos} & $1.873 \times 10^{-2}$ & $1.639 \times 10^{-2}$ \\
\texttt{small\_rel\_dsq} & $1.960 \times 10^{-2}$ & $1.172 \times 10^{-2}$ \\
\texttt{small\_rel\_idc} & $2.594 \times 10^{-2}$ & $2.073 \times 10^{-2}$ \\
\end{tabular}
\caption{The smaller models.}
\label{fig:results-small}

\end{figure}

The \texttt{small} GKAN model without any computed features performs about 1 order of magnitude worse compared to the baseline GNN. This alone is quite impressive when considering the difference in model size.

Furthermore, by introducing simple computed features such as $\Delta x$ and $\Delta y$ (which would not be unreasonable to include, even without \textit{a priori} knowledge of the desired force law), we observe significantly improved performance on par with the baseline---with the additional benefit of a 99.5\% reduction in model size. This suggests that the GKAN could be a highly efficient architecture for domain-informed use cases.

\ag{Include loss plots and rollouts here?}

\section{Discussion (REAL)}
\subsection{Regularization}
None of the different regularizations used resulted in networks sparse enough to be interpretable through visual inspection. Even using multipliers of 10 or 100 for the overall regularization only produced marginally sparser GKANs compared to the no-regularization control, while being unable to learn due to the regularization loss overpowering prediction loss. It was expected that the L1-norm and entropy regularizations alone would be insufficient, though even combined, as outlined in \cite{liu2024kankolmogorovarnoldnetworks}, results are disappointing. This may be explained by:
\begin{enumerate}
    \item The gravitational $n$-body problem requiring dense networks to accurately represent the dynamic system. Though  \cite{liu2024kan20kolmogorovarnoldnetworks} provides an example of a symbolic KAN modeling gravitational force between two bodies in a small, interpretable network, indicating that the gravitational forces do not require dense networks. \ek{I HATE THIS I HATE THIS THIS}
    \item Constant regularization does not allow for balance between learning and sparsity. The experiments did not increase regularization magnitudes over the learning period, which may have done a better job of letting the model optimize for accurate predictions first, then learn a sparse representation after.
\end{enumerate}

Why does all ts suck so much ass? And why no sparse?
\begin{itemize}
    \item maybe gravity just requires dense networks to model without using feature engineering; having to learn relative positions, multiplication and stuff might just be too hard for our little guys
    \item maybe the nature of KANs just doesn't lend itself to learning networks
    \item maybe if we had a bajillion more epochs it would actually fix itself

    \hs{igb...}
    \hs{maybe think about potential improvements?}
\end{itemize}

\subsection{How small KAN we get?}

\ag{Could do interpretability analysis. Pass synthetic inputs through the network to see the relationships empirically (sanity check). Then try symbolic regression per-layer and analytically composing them to see if it recovers the force law?}

\ag{The following is outdated.} Looking at the visualizations, it seems that the models without any augmentation did not learn as effectively. The message and node update networks are both quite dense, which suggests that they did not learn to ignore distractors such as velocity. By augmenting the dataset with \texttt{rel\_pos}, we introduce an inductive bias that is apparently sufficient to ``push'' the model in the right direction, and interestingly, we can start to see sparsity in the networks. Refer to the visualizations for \texttt{minimal} and \texttt{minimal\_rel\_pos} to see this effect in action. \hs{link to KAN 2.0 paper}

I have no idea what mathematical relationships it actually learned. \texttt{/shrug} \hs{show image}

\hs{plots and rollouts? maybe keep in appendix}

\subsection{How small KAN we get?}

\ag{Could do interpretability analysis. Pass synthetic inputs through the network to see the relationships empirically (sanity check). Then try symbolic regression per-layer and analytically composing them to see if it recovers the force law?}

\ag{The following is outdated.} Looking at the visualizations, it seems that the models without any augmentation did not learn as effectively. The message and node update networks are both quite dense, which suggests that they did not learn to ignore distractors such as velocity. By augmenting the dataset with \texttt{rel\_pos}, we introduce an inductive bias that is apparently sufficient to ``push'' the model in the right direction, and interestingly, we can start to see sparsity in the networks. Refer to the visualizations for \texttt{minimal} and \texttt{minimal\_rel\_pos} to see this effect in action. \hs{link to KAN 2.0 paper}

I have no idea what mathematical relationships it actually learned. \texttt{/shrug} \hs{show image}

\hs{plots and rollouts? maybe keep in appendix}

\subsection{KAN't Optimize}

% Overall, our results suggest that the main bottleneck for GKANs in this setting may not be representational power, but trainability. Better optimization and stronger physical inductive biases may be necessary before the interpretability advantages of KAN-based graph models can be fully realized.

% discuss difficulties with learning usng GKANs (tests with differently conditioned problems, min separation distance in data, poor loss landscape, L2 vs L1 losses)

% dicuss things we tried (L-BFGS, Adam, hybrid opitmization, alternating optimization (nothing works)). Why is this different from KAN paper (they use L-BFGS)

% discuss challenges of using message passing architecture (could this be adding irregularity to loss landscape? preventing convergence and proper learning)

% discuss infeasiblility for certain asymptotic functions for spline edges, discuss potential tricks with grid updates etc. 

discuss potential for physically informed loss function (e.g., incorporating conservation laws or physical constraints into the loss function, hamiltonian loss functions, etc). THis may partially explain poor learning results over all

discuss future improvements like pruning low-weighted edges every few epochs

\newpage

\printbibliography

\appendix

\hs{we need to add github links, rollout, charts, images}

\section{Supplementary material} \label{app:info}

This stuff doesn't count towards your page limit.

\newpage
\section{What goes in the report}
Part of your grade is based on the contribution of your project:
did you do something meaningful in your project, that adds something new to the world?
Again, we don't expect you to necessarily make a high-impact totally novel result,
and it's fine if you end up with a negative result,
but you should have something here.
You should make sure that the contribution of your project is clear in the introduction of your paper,
and that the rest of the writeup clearly demonstrates that you actually made this contribution.

A large portion of your grade is based on writing a clear and structured paper that addresses the relevant questions that would be asked of a generic scientific/engineering publication.
Your writeup should look more or less like a scientific paper as published at, say, NeurIPS,
or a NeurIPS workshop.
To achieve that, most papers should use something like the following traditional outline:
\begin{enumerate}
\item Introduction: Clearly state the problem being addressed. \textbf{Explain why it is an important problem to work on}. At a high level, briefly summarize what the limitations of existing approaches that your work will be addressing, and what the contribution of your project is.
\item Related Work: Identify at least three publications on related topics; usually these will be papers that have worked on slightly different problems or papers that have proposed an approach to your problem that is not fully satisfactory. For each paper, briefly say either how the problem addressed is related and/or different (if they address a different problem) or why it doesn't solve the problem you are working on (if it addresses the same problem). Sometimes it makes more sense to put this later in the paper, but usually it's here.
\item Description and justification of what you did, divided up (possibly in multiple sections and/or sub-sections). There's a lot of flexibility here, and it will depend on the type of project you are doing. For example, if you're applying standard machine learning methods to a new dataset or doing a Kaggle competition, you could have one (sub-)section describing the dataset and why you think machine learning could help, and one (sub-)section stating the methods you will try and why you think these are appropriate methods (you don't necessarily have to go into detail describing the methods). If you're extending an existing technique, you could have one sub-section describing the existing technique, and one sub-section for each of the extensions you explored. If your project has a theoretical component, you might have one sub-section discussing the assumptions, one sub-section describing the results, and one sub-section describing implications.
\item Experiments and/or analysis (if you have an experimental component): Describe each experiment that you did. Say what each experiment is trying to test. Ideally, each experiment should only try to test one thing and you should control for as many other factors as possible. Subsequently, summarize the result of your experiment, in both the text and in a nice visual form such as a figure; \textbf{a table with a huge list of numbers is often not the best way to summarize information}, but sometimes a well-designed table (especially with some colour-coding) is still your best bet.
\item Discussion and future work: State the main conclusions you obtained from this course project. List at least one strength and one weakness of your contribution. Briefly describe what you might do with more time.
\end{enumerate}

An example outline for a perspective paper might be something like
\begin{enumerate}
\item Introduction: Clearly state the problem being addressed. Explain why it is an important problem. At a high level, briefly summarize the history of the works that will be discussed in the project.
\item Review: Go through the different works in some logical order, such as chronologically or by going from simple to complex models. Don't just list the methods, but say how they relate to each other (going through the strengths/weaknesses of the different methods, both in comparison to each other and compared to an ideal method that solves the problem).
\item Discussion: Discuss the trends that have occurred over time. Speculate about where the next steps in the trend could lead. Point out issues that are not properly addressed by existing methods. State some interesting directions to explore, or opportunities to use existing tools in new applications.
\end{enumerate}

You definitely don't have to stick to these exact structures; many of my papers don't.
But if you're not using this format (\dots or if you are), you should make sure what you're doing makes sense,
and answers the important questions about your project.

\paragraph{Writing style}
Mark has some advice for writing here, most of which I agree with:
\url{https://www.cs.ubc.ca/~schmidtm/Courses/Notes/writing.pdf}.
(He has a bunch of minor grammatical errors in that document, though. :/)

You're not going to lose points for making some minor grammatical mistakes or anything like that
-- the published literature is full of them!
You will lose points if there are grammatical errors to the level of making your work hard to understand.
More importantly, you will also lose points if your writing doesn't convey what you did
in a clear, easy-to-understand way.
Note that this involves a lot more than just having a good grasp of written English.


\section{How to format stuff}
\subsection{Citations}

Something a lot of people don't know at first:
use \verb+\citet{ref}+ (or \verb+\textcite+)
if the citation plays a grammatical role in the sentence,
e.g.\ ``\citet{transformers} demonstrated that \dots.''
Use \verb+\citep{ref}+ (or \verb+\parencite+)
if it doesn't, e.g.\ ``Machine learning is fun \citep{ml-is-fun}.''


\subsection{Figures}

See \cref{fig:sample} for how to include a figure.

\begin{figure}
  \centering
  % The next line would normally be \includegraphics instead.
  \fbox{\rule[-.5cm]{0cm}{4cm} hi =) \rule[-.5cm]{4cm}{0cm}}
  \caption{Sample figure caption.}
  \label{fig:sample}
\end{figure}


\subsection{Tables}

Table captions go \emph{above} the table,
because that's the usual style, idk.
There's an example in \cref{tab:sample}.
Avoid having too many lines, and in particular don't just put them between every table cell; put them where they help readability.

\begin{table}
  \caption{Sample table title}
  \label{tab:sample}
  \centering
  \begin{tabular}{lll}
    \toprule
    \multicolumn{2}{c}{Part}                   \\
    \cmidrule(r){1-2}
    Name     & Description     & Size ($\mu$m) \\
    \midrule
    Dendrite & Input terminal  & $\sim$100     \\
    Axon     & Output terminal & $\sim$10      \\
    Soma     & Cell body       & up to $10^6$  \\
    \bottomrule
  \end{tabular}
\end{table}

\subsection{Math}
Using display math in bare TeX commands, like \texttt{\$\$f(x)\$\$}, will mess up the line numbers. (Inline math like \texttt{\$f(x)\$} is fine.)
Instead use \texttt{\textbackslash [ f(x) \textbackslash ]} or \texttt{\textbackslash begin\{equation*\} f(x) \textbackslash end\{equation*\}},
like this:
\[ \sum_i a_i = 2 .\]
(You shouldn't use \texttt{\$\$} anyway; see \href{https://tex.stackexchange.com/questions/503/why-is-preferable-to}{this discussion} and \href{https://tex.stackexchange.com/questions/40492/what-are-the-differences-between-align-equation-and-displaymath}{this one} for more.)


\subsection{Supplementary Material}
You can include extra information in appendices, like \cref{app:info}.
You don't have to if you don't want to.

Please \emph{don't} include code in your writeups, unless you did something particularly cool and want to \emph{briefly} describe the way something works as a contribution (this should be relatively unusual).
Please instead link to it somewhere, preferably a GitHub repo or similar,
especially if the code is a significant contribution of your project.


\section{Discussion}
In the end, you should give us an A++++.


\begin{ack}
    You can acknowledge useful discussion with other people who aren't coauthors here (or leave the section out).
    Typically you'd also put funding here, acknowledgements for computing clusters, etc.
\end{ack}


% \printbibliography
\printbibliography[title={References}]


%%%%%%%%%%%%%%%%%%%%%%%%%%%%%%%%%%%%%%%%%%%%%%%%%%%%%%%%%%%%



\end{document}